\section*{Abstract}

The rapid development of sports analytics over the past two decades has yielded increasingly sophisticated tools for the assessment of individual athletic performance. 
From traditional box-score statistics to granular tracking-derived indicators, the field has devoted considerable effort to quantifying and comparing players at the individual level. 
However, in team sports, it is ultimately collective units, and not isolated individuals, that compete and win. 
How, then, should individual performance data be leveraged to inform the construction of a competitive team? 
This question, central to roster management, remains rarely addressed in the literature.
In this paper, we propose a two-step, data-driven framework for optimal lineup building in basketball, illustrated on five National Basketball Association (NBA) seasons. 
First, we apply and benchmark several unsupervised clustering algorithms on a rich dataset of player-season observations combining box-score and tracking statistics, with the objective of identifying robust, interpretable player archetypes that transcend the individual level. 
Second, we leverage these archetypes to evaluate lineup compositions extracted from play-by-play data, characterizing each combination of profiles by its associated plus-minus per 48 minutes differential, a direct measure of a lineup's net contribution to the game outcome, aggregated across a statistically sufficient number of appearances.
Our results confirm that, while "superteams" tend to perform well, lineup efficiency is driven not by the mere accumulation of individually dominant players, but by the complementarity of the profiles they embody. 
These findings carry concrete implications for roster construction and strategic decision-making, and open promising avenues for future work incorporating financial constraints and sequential team-building frameworks.